\documentclass[a4paper,10pt]{article}
\usepackage[top=14mm,bottom=16mm,left=15mm,right=15mm]{geometry}
\usepackage{fontspec}
\usepackage{xeCJK}
\setmainfont{Times New Roman}
\setCJKmainfont{Yu Gothic}
\usepackage{tabularx}
\setlength{\parindent}{0pt}
\newcommand{\answerrow}[2]{\textbf{#1}\quad #2\\[-1mm]\fbox{\parbox[t][38mm][t]{0.96\linewidth}{\hspace{1mm}}}\par\vspace{3mm}}
\begin{document}
\begin{center}
{\LARGE\bfseries 車両同定MLE 完全手計算体験書}\\[2mm]
{\Large 解答用紙}\\[5mm]
\begin{tabularx}{\linewidth}{|p{25mm}|X|p{28mm}|}\hline
氏名&&採点\\\hline 実施日&&確認者\\\hline
\end{tabularx}
\end{center}
\vspace{4mm}
\answerrow{M1}{Gaussian NLL}
\answerrow{M2}{Huber}
\answerrow{M3}{Hampel}
\answerrow{M4}{map shape}
\answerrow{M5}{PV gain}
\answerrow{M5b}{DNI/DHIから停止POA}
\answerrow{M5c}{時間平均日射と瞬時日射}
\answerrow{M5d}{DC bus発電ゲインと二重補正}
\answerrow{M5e}{25直列pack電圧上限}
\answerrow{M6}{motion power}
\answerrow{M6c}{inertial energy}
\answerrow{M6b}{70 km/h総負荷検算}
\answerrow{M7}{識別可能性}
\answerrow{M8}{battery IV/SoC}
\answerrow{M9}{voltage residual}
\answerrow{M10}{terminal OCV}
\answerrow{M11}{joint objective}
\answerrow{M12}{有限差分}
\answerrow{M13}{multi-start}
\answerrow{M14}{map採否}
\answerrow{M15}{canonical profile}
\answerrow{M16}{summary検査}
\answerrow{M17}{小RMSEの落とし穴}
\answerrow{M18}{全工程}
\end{document}
